Un cylindre, de volume $V=\qty{2}{\cubic\m}$, contient un gaz de dioxyde
de carbone \ce{CO2} à la température $\theta=\qty{20}{\degreeCelsius}$
et sous une pression de $P_1=\qty{1013}{\hPa}$. A température
constante, on ajoute une quantité de gaz de dioxygène \ce{O2} et la
pression du mélange des gaz devient $P_2=\qty{1040}{\hPa}$.
\begin{enumerate}
    \item Calculer $n_1$ la quantité de matière de \ce{CO2} qui se
          trouve dans le cylindre.
    \item Calculer la masse $m$ du mélange gazeux qui se trouve dans
          le cylindre.
\end{enumerate}

\begin{donnees}
    \begin{itemize}
        \item $M(S)=\qty{32}{\g\per\mol}$, $M(O)=\qty{16}{\g\per\mol}$,
              $M(C)=\qty{12}{\g\per\mol}$, $M(H)=\qty{1}{\g\per\mol}$
        \item La constant d'Avogadro~: $N_A=\qty{6.02e23}{\per\mol}$
        \item La masse volumique de l'eau~: $\rho_e=\qty{1}{\g\per\mL}$
        \item La constante des gaz parfait~:
              $R=\qty{8.31}{\Pa\cubic\m\per\mol\per\K}$
    \end{itemize}
\end{donnees}
